\textbf{COLUNI 2022 -- Questão 01}\par\medskip Leia o texto a seguir para as questões 01 e 02.

\par

\begin{center}\textbf{História em quadrinhos}\end{center}

As histórias em quadrinhos (HQs) são amplamente conhecidas, principalmente pelos clássicos de heróis e vilões. No entanto, o gênero não se restringe a essa possibilidade, e, na realidade, é utilizado como recurso narrativo para diversas histórias, que podem ser ficcionais ou representativas de fatos históricos. Caracterizam-se pela predominância da linguagem visual como recurso comunicativo, no qual as cenas individuais são postas em sequência para construir-se um movimento narrativo. Nascidas e expandidas pelo surgimento da imprensa e do jornal, as histórias em quadrinhos conquistaram sua autonomia e ampliaram seu alcance, adentrando, inclusive, o cinema. As diversidades estéticas e temáticas solidificaram a força do gênero. [...]

\par

\textbf{O que é história em quadrinhos?}

História em quadrinhos (HQ) é um texto de linguagem mista, no qual predomina o uso da linguagem visual como recurso narrativo. Além disso, esse gênero define-se por sua estrutura em pequenos quadros [...] as imagens e possíveis falas expressam-se em uma relação de tempo-espaço no texto. Comumente é conhecida ou como gênero lúdico, voltado principalmente ao público infantil, ou como as famosas histórias de super-heróis, que, pelo grande sucesso que tiveram e ainda têm no grande público, confundem-se com o gênero. Entretanto, existem diversos temas, estéticas, linguagens e universos trabalhados por histórias em quadrinhos e que nem sempre são reconhecidos, como as adaptações de textos literários, os quadrinhos de cunho mais filosófico etc. Existem diversos tipos de quadrinhos, com estéticas e histórias próprias, como os de terror, romance, ficção científica, os que fazem releituras de outros textos, os de gênero crime/policial, infantil, juvenil, os filosóficos, históricos etc. [...] Os quadrinhos, por sua predominância textual na imagem, retratam os movimentos e o passar do tempo por meio da sequencialidade de imagens que representam momentos estáticos, mas que, vistos em relação aos outros quadrinhos, contam uma história. Existem muitas HQs que só utilizam a linguagem visual e exploram as aproximações e os distanciamentos, as perspectivas, os exageros, as cores e outros elementos, para representar ou expressar comportamentos, que são compreendidos pelo leitor. [...]

\par

Os quadrinhos, como os conhecemos hoje, nasceram com o uso e expansão da imprensa, pois, pelos jornais, as primeiras tirinhas foram publicadas e consumidas pelos leitores. Esse fato garantiu ao gênero uma inserção no consumo cultural, mas, por outro lado, limitava suas publicações, pois os autores dependiam do jornal. Por volta da década de 1930, as tirinhas já publicadas em jornais passaram a ser republicadas em revistas independentes. Como o jornal geralmente era descartado após a leitura, essas revistas cumpriam a função de manter as tirinhas à disposição do público. [...] Com o sucesso das primeiras revistas em quadrinhos, o mercado para esse gênero cresce e novas histórias são criadas e publicadas, novos estilos consolidam-se, além de temas e representações políticos que passaram a ser debatidos nelas. O sucesso dessas revistas alcançou também o mercado cinematográfico, e diversas histórias passaram a ter a sua versão em filme. [...]

\par

Disponível em: https://mundoeducacao.uol.com.br/literatura/historia-historia-quadrinhos.htm. Acesso em: 05 set. 2021 (Adaptado).

\par

Analise as afirmativas abaixo em relação ao texto “História em quadrinhos”.

\par

I. No texto, ressalta-se a importância do gênero HQs no contexto da comunicação.\par
II. O tempo e o espaço nas histórias em quadrinhos, na maioria das vezes, não são mencionados.\par
III. História em quadrinhos pressupõe, em geral, uma mescla de linguagens verbal e não verbal.\par
IV. Os quadrinhos são sempre elaborados a partir de adaptações de textos clássicos ou literários.

\par

Está CORRETO o que se afirma apenas em:

\par

a) I e II.\par
b) I e III.\par
c) II e IV.\par
d) III e IV.
